排序和区间深度题卡 不再把所有题目堆成一个清单，而是按题型分成章内大节。每一节先讲分流规则，再列代表案例，最后逐题做单点分析。

\section{区间、排序与合并：题型分流与案例}
这一节先把同一题型的常见场景拆开讲清楚，再进入具体题目。学习顺序是：先写暴力基线，定位最贵的重复工作；再根据题目条件选择能维护不变量的数据结构；最后用多个案例比较边界、反例和变体。

\begin{principlebox}{图示：区间、排序与合并推导链}
\begin{tabular}{@{}p{0.18\linewidth}p{0.74\linewidth}@{}}
暴力基线 & 枚举候选、完整模拟或逐个检查，先保证覆盖全部可能答案。\\
瓶颈定位 & 瓶颈是没有利用排序后相邻关系。区间按起点或终点排序后，很多远离的区间不可能再影响当前决策。\\
结构选择 & 合并按起点排序，调度按终点排序，覆盖题常用起点升序且终点降序。扫描时维护当前最远右端、结果尾区间或已选区间终点。\\
不变量 & 结果数组中的区间已经互不重叠且有序；当前扫描区间只可能影响结果最后一个区间或当前边界。\\
代码落地 & 比较器必须处理相同起点，否则覆盖和去重会错。区间端点可能溢出时用更大整数。闭区间和半开区间决定是否用小于等于。\\
\end{tabular}
\end{principlebox}

\begin{principlebox}{题型分流：区间、排序与合并}
\textbf{合并和插入。}
判断条件：区间按起点排序后，重叠区间连续出现。处理方式：维护结果尾区间或新区间边界。代表案例：56 合并区间、57 插入区间。风险点：闭区间和半开区间端点相接处理不同。

\textbf{覆盖和删除。}
判断条件：判断一个区间是否被之前更大的右端覆盖。处理方式：同起点终点降序，扫描维护最大右端。代表案例：1288 删除被覆盖区间、435 无重叠区间。风险点：同起点排序方向是关键。

\textbf{调度和最少资源。}
判断条件：多个区间竞争时间资源，问最多保留或最少会议室。处理方式：按终点贪心或按起点扫描堆维护当前结束时间。代表案例：252 会议室、253 会议室 II、452 引爆气球。风险点：端点相等是否冲突由题意决定。

\textbf{动态区间。}
判断条件：区间逐个插入，必须立即判断冲突或维护覆盖。处理方式：有序表查前驱后继，或扫描线维护事件。代表案例：729 我的日程安排表 I、731 我的日程安排表 II。风险点：哈希表无法回答前驱后继。

\end{principlebox}

本节案例覆盖 LeetCode 56 合并区间、LeetCode 57 插入区间、LeetCode 252 会议室、LeetCode 253 会议室 II、LeetCode 352 将数据流变为多个不相交区间、LeetCode 436 寻找右区间、LeetCode 729 我的日程安排表 I、LeetCode 731 我的日程安排表 II、LeetCode 986 区间列表的交集、LeetCode 1094 拼车、LeetCode 1109 航班预订统计、LeetCode 1229 安排会议日程、LeetCode 1288 删除被覆盖区间。阅读时不要按题号背答案，而要先判断题面落在哪个分流场景：查询目标是什么，输入约束给了什么单调性或等价关系，暴力慢在哪里，优化结构保存了什么状态。

\subsection{代表案例与单点分析}

\subsubsection{LeetCode 56 合并区间}

\textbf{题目模型。} 排序后可能与当前区间相交的只会是结果尾部。

\textbf{推导重点。} 按起点排序，扫描时维护结果最后区间的右端点。

\textbf{边界。} 端点相接、完全覆盖、无重叠、空列表、输入未排序。

\textbf{变体追问。} 若区间是半开区间，端点相等不一定合并。

\subsubsection{LeetCode 57 插入区间}

\textbf{题目模型。} 新区间只会影响与它相交的一段连续区间。

\textbf{推导重点。} 先追加左侧完全不相交区间，再合并重叠段，最后追加右侧。

\textbf{边界。} 新区间在最前、最后、覆盖全部、刚好相邻、原列表为空。

\textbf{变体追问。} 若输入不是有序且不重叠，必须先排序或换模型。

\subsubsection{LeetCode 252 会议室}

\textbf{题目模型。} 每个会议占用一个时间区间，任意重叠都表示不能全部参加。

\textbf{推导重点。} 按开始时间排序，只需检查相邻会议是否冲突。

\textbf{边界。} 端点相等是否冲突、空列表、单会议、完全重叠。

\textbf{变体追问。} 若问最少会议室，需要维护同时进行的会议数量。

\subsubsection{LeetCode 253 会议室 II}

\textbf{题目模型。} 同时进行的会议数量决定最少房间数。

\textbf{推导重点。} 按开始时间扫描，用小顶堆维护当前会议的最早结束时间。

\textbf{边界。} 会议端点相接、多个同一时间开始、空列表、堆顶释放条件。

\textbf{变体追问。} 也可以把开始和结束分别排序，用双指针统计并发数。

\subsubsection{LeetCode 352 将数据流变为多个不相交区间}

\textbf{题目模型。} 数字逐个到达，需要动态维护已覆盖的有序不相交区间。

\textbf{推导重点。} 用有序表找到新值附近的前驱和后继，判断是否连接、合并或新建区间。

\textbf{边界。} 重复插入、连接左右两个区间、插入最小最大值、相邻端点。

\textbf{变体追问。} 这类动态区间不能用普通数组扫描长期维护。

\subsubsection{LeetCode 436 寻找右区间}

\textbf{题目模型。} 每个区间需要寻找起点不小于当前终点的最小起点区间。

\textbf{推导重点。} 把所有起点和原下标排序，对每个终点做 lower bound。

\textbf{边界。} 不存在右区间、起点相同、负端点、单区间。

\textbf{变体追问。} 如果区间动态插入，则需要有序表而不是一次排序数组。

\subsubsection{LeetCode 729 我的日程安排表 I}

\textbf{题目模型。} 每次插入新区间时，需要判断它是否和已有区间重叠。

\textbf{推导重点。} 有序表按起点存区间，只检查新区间的前驱和后继。

\textbf{边界。} 端点相接、插入最前最后、完全覆盖已有区间、动态多次插入。

\textbf{变体追问。} 哈希表无法回答前驱后继，动态区间要用有序结构。

\subsubsection{LeetCode 731 我的日程安排表 II}

\textbf{题目模型。} 允许双重预订但不允许三重预订，插入会改变区间重叠层数。

\textbf{推导重点。} 保存已有预订和已经双重预订的区间；新预订若碰到双重区间则失败，否则记录新产生的重叠。

\textbf{边界。} 端点相接、多个局部重叠、完全覆盖、回滚失败插入。

\textbf{变体追问。} 扫描线差分也能做，但要处理失败时恢复状态。

\subsubsection{LeetCode 986 区间列表的交集}

\textbf{题目模型。} 两个有序且不重叠的区间列表，交集只可能来自当前两个区间。

\textbf{推导重点。} 每次比较两个当前区间的重叠部分，然后推进右端较小的区间。

\textbf{边界。} 空列表、端点相接、一个区间覆盖多个区间、无交集。

\textbf{变体追问。} 这是区间双指针，不需要把两个列表合并排序。

\subsubsection{LeetCode 1094 拼车}

\textbf{题目模型。} 乘客上下车事件构成沿路线位置变化的载客量。

\textbf{推导重点。} 用差分数组或扫描线记录每个位置人数变化，累计人数不能超过容量。

\textbf{边界。} 同站上下车、容量刚好、站点范围、输入未排序。

\textbf{变体追问。} 航班预订统计也是差分思想，只是只需要最终每段累加结果。

\subsubsection{LeetCode 1109 航班预订统计}

\textbf{题目模型。} 每条预订给一个连续航班区间增加座位数。

\textbf{推导重点。} 差分数组在区间起点加人数，终点后一位减人数，最后前缀还原每航班总数。

\textbf{边界。} 区间到最后一个航班、单点区间、多条重叠预订、下标从一开始。

\textbf{变体追问。} 这是静态区间加法，不需要线段树。

\subsubsection{LeetCode 1229 安排会议日程}

\textbf{题目模型。} 两个人的可用时间段都有序，目标是找最早长度足够的交集。

\textbf{推导重点。} 双指针比较当前两个区间交集，若长度不足就推进结束更早的一侧。

\textbf{边界。} 无可行时间、刚好等于 duration、区间端点相接、输入需要先排序。

\textbf{变体追问。} 它和区间列表交集相同，但额外带最早可行和长度约束。

\subsubsection{LeetCode 1288 删除被覆盖区间}

\textbf{题目模型。} 被覆盖要求另一区间起点不晚且终点不早。

\textbf{推导重点。} 按起点升序、终点降序排序，扫描维护最大右端。

\textbf{边界。} 同起点区间、完全覆盖、无覆盖、端点相等。

\textbf{变体追问。} 终点降序是关键，否则同起点小区间会先出现造成误判。
